% ==============================================================================
% =================================== Einleitung ================================
% ==============================================================================

%\chapter{Einleitung}
%\label{chap:einleitung}
%\thispagestyle{fancy}

\section{Einleitung}

Die fortschreitende Digitalisierung und die Verbreitung mobiler Endgeräte generieren kontinuierlich gesundheitsbezogene Sensordaten, die für die personalisierte Medizin genutzt werden können, jedoch eine adäquate Aufbereitung erfordern \cite[4]{2}\cite[1]{3}. Ziel dieser Arbeit ist es, ein Visualisierungskonzept für ein klinisches Frühwarnsystem zu entwickeln, das Wearable-Daten (Herzfrequenz, Blutdruck, SpO₂, Bewegungs-, Schlafdaten, Patientengruppe) sowohl für die explorative Ärzte-Analyse als auch für die Patientenkommunikation aufbereitet \cite[8]{2}.

Dazu werden die relevanten Datenattribute theoretisch eingeordnet, zielgruppenspezifische Diagrammtypen ausgewählt und ethische Aspekte beleuchtet. Die praktische Umsetzung erfolgt mit der Programmiersprache R.

Die Arbeit gliedert sich wie folgt: Kapitel 2 widmet sich den theoretischen Grundlagen der Datenvisualisierung und der E-Notation. Kapitel 3 behandelt die zielgruppengerechte Gestaltung von Visualisierungen sowie ethische und wahrnehmungspsychologische Aspekte. Kapitel 4 stellt die praktische Umsetzung in R dar, gefolgt von einer abschließenden Zusammenfassung und Diskussion der Ergebnisse.

Die Zitierweise folgt dem numerischen System. Die Quellen werden mit den Schlüsseln \cite{1}, \cite{2} und \cite{3} referenziert.

% ==============================================================================
% =================================== Einleitung ================================
% ==============================================================================